\section{Conclusion}

This paper presents a secure, explainable, blockchain-audited overlay for
inter-agency police record access governance. The main idea is simple:
sensitive record sharing needs contextual access control, tamper-evident audit,
and explanations that can be reviewed later. None of these parts is sufficient
alone.

The proposed work contributes a layered access-control service for
allow/deny/escalate decisions, a permissioned blockchain-style audit layer for
commitments and decision evidence, and an explainable decision service that
records why each access request was allowed, denied, or escalated. The
benchmark also studies a validly re-signed compromised-signer attack, where a
log can appear valid even after policy output is corrupted.

The work remains intentionally bounded. It is a synthetic benchmark and
research prototype, not an operational police deployment. The results support
a layered interpretation: blockchain-style audit improves tamper evidence for
recorded events, contextual policy checking improves decision accountability,
and explainable drift/audit artifacts help review failure modes that are not
visible to a single mechanism alone. Future work should validate policies with
domain experts, improve the natural-language explanation renderer, strengthen
privacy analysis, and test the design under officially approved operational
constraints.
